% Generated from locale/tr/content/lambda-calculus/syntax/de-bruijn.tex; do not hand-edit.


\olfileid[tr]{lam}{syn}{deb}

\olsection{De Bruijn İndisi}

$\alpha$-denkliği son derece doğaldır; çünkü $\alpha$-denk terimler ``aynı
anlama gelir.'' Aslında birbirine $\alpha$-dönüşebilen terimleri ayırt etmeyen
bir lambda terimleri sözdizimi vermek mümkündür. Bunların en iyi bilineni,
değişkenlerin yerine \emph{De Bruijn indislerini} kullanır.

$\lambd[x][M]$ yazdığımızda $x$'in fonksiyonun parametresi olduğunu açıkça
belirtir ve $M$ içinde bu parametreye $x$ ile gönderme yaparız. De Bruijn
gösteriminde ise parametrelerin adı yoktur; fonksiyon gövdesinde bir
parametreye yapılan gönderme, gönderme ile onu bağlayan soyutlama arasındaki
soyutlama katmanlarının sayısıyla gösterilir. Örneğin
$\lambd[x][\lambd[y][y x]]$ terimini ele alalım: dıştaki soyutlama $x$
değişkenini, içteki soyutlama $y$ değişkenini bağlar. $y x$ alt terimi
içteki soyutlamanın kapsamındadır. $y$ ile onu bağlayan $\lambd[y]$ arasında
başka soyutlama yoktur; $x$ ile onu bağlayan $\lambd[x]$ arasında ise bir
soyutlama vardır. Bu nedenle $y x$ için $0\,1$, bütün terim için de
$\lambd[][\lambd[][01]]$ yazarız.

\begin{defn}
  De Bruijn terimleri tümevarımlı olarak şöyle tanımlanır:
  \begin{enumerate}
  \item Herhangi bir doğal sayı $n$;
  \item $P$ ve $Q$ birer De Bruijn terimi olmak üzere $PQ$;
  \item $N$ bir De Bruijn terimi olmak üzere $\lambd[][N]$.
  \end{enumerate}
\end{defn}

Alışılmış lambda terimlerinden De Bruijn indisli terimlere biçimselleştirilmiş
çeviri şöyledir:
\begin{defn}
  \begin{align*}
    F_\Gamma(x) &= \Gamma(x) \\
    F_\Gamma(PQ) &= F_\Gamma(P)F_\Gamma(Q) \\
    F_\Gamma(\lambd[x][N]) &= \lambd[][F_{x,\Gamma}(N)]
  \end{align*}
  Burada $\Gamma$, sıfırdan başlayarak indislenmiş bir değişkenler listesidir;
  $\Gamma(x)$, $x$ değişkeninin $\Gamma$ içindeki konumunu gösterir. Örneğin
  $\Gamma$ listesi $x,y,z$ ise $\Gamma(x)=0$ ve $\Gamma(z)=2$ olur.

  $x,\Gamma$, $x$'in $\Gamma$ listesinin başına eklenmesiyle elde edilen
  listeyi gösterir. Örneğin yukarıdaki $\Gamma$ için $w,\Gamma$ listesi
  $w,x,y,z$'dir.
\end{defn}

Bir De Bruijn teriminden standart lambda terimini geri elde etme işlemi
şöyledir:

\begin{defn}
  \begin{align*}
    G_\Gamma(n) &= \Gamma[n] \\
    G_\Gamma(PQ) &= G_\Gamma(P) G_\Gamma(Q) \\
    G_\Gamma(\lambd[][N]) &= \lambd[x][G_{x,\Gamma}(N)]
  \end{align*}
  Burada $\Gamma$ yine sıfırdan başlayarak indislenmiş bir değişkenler
  listesidir; $\Gamma[n]$, $n$ konumundaki değişkeni gösterir. Örneğin
  $\Gamma$ listesi $x,y,z$ ise $\Gamma[1]=y$ olur.

  Son eşitlikteki $x$, $\Gamma$ içinde bulunmayan herhangi bir değişken olarak
  seçilir.
\end{defn}

Şimdi bazı sonuçları ispatsız veriyoruz:

\begin{prop}
  $M \aconvone M'$ ve $\Gamma$, $\FV{M}$ kümesini içeren herhangi bir
  listeyse $F_\Gamma(M) \eqs F_\Gamma(M')$ olur.
\end{prop}

